%% kapa: exam=1 round=21 year=2010 period=2 form=A subject=appraisal_law subject_ko=부동산관계법규 seq=18 number=58
%% source: Q-Net 감정평가사 기출문제 (cst00309) artlSeq=5024181 file=extracted/1008/5024181/2010년 제21회 감정평가사 1차시험 기출문제/2교시.hwp
%% conversion: kordoc 4.24.51 -> tools/gen1.py (auto); review=proofread-gpt-sol; flags=none
\begin{question}{58}
부동산 가격공시 및 감정평가에 관한 법령상 표준지공시지가에 관한 설명으로 옳지 않은 것은?
\begin{choices}
  \item 토지거래 또는 임대차에 대한 지식부족으로 인하여 이루어진 인근 유사 토지의 거래가격 또는 임대료는 표준지의 적정가격의 조사․평가 기준이 될 수 없다.
  \item 표준지공시지가의 공시사항에는 표준지의 단위면적당 가격, 표준지 및 주변토지의 이용상황, 도로․교통상황, 지세 등이 포함되어야 한다.
  \item 표준지공시지가에 대한 이의신청은 서면으로 국토해양부장관에게 하되 토지소유자, 토지이용자 등 법률상․사실상 이해관계인이 할 수 있다.
  \item 표준지공시지가는 일반적 토지거래의 지표가 되며, 감정평가업자가 개별적으로 토지를 감정평가하는 경우 그 기준이 된다.
  \item 표준지공시지가는「농어촌정비법」에 따른 농업생산기반 정비사업을 위한 환지․체비지의 매각 또는 환지신청에 적용된다.
\end{choices}
\end{question}
\answer{3}
